\section{Entorno de trabajo en Python sobre MSTM}
\label{sec:python_wrapper}

El código MSTM~\cite{Mackowski_code} se controla mediante un fichero
de texto plano (\texttt{mstm.inp}) con sintaxis numérica de Fortran,
devuelve los resultados en ficheros \texttt{.dat} cuya estructura
depende de las opciones activadas, y trabaja exclusivamente en unidades
adimensionales escaladas por el número de onda. Para automatizar el
ciclo completo de una simulación ---configuración, conversión de
unidades, ejecución, lectura de salidas y visualización--- se ha
desarrollado un \emph{wrapper} en Python. El código completo
está disponible en GitHub:
\url{https://github.com/GorkaLarrazabalEhu/mie-scattering-project}.

El entorno se diseñó para poder explorar distintas configuraciones y
observables, no solo los que se emplean en este trabajo. Por eso
incluye más opciones de las que luego aparecen en los resultados; en
esta sección se describe el conjunto general y, cuando procede, se
indica qué parte se usa en cada caso.

El wrapper se organiza en una interfaz gráfica, un módulo de utilidades
físicas y un conjunto de scripts de representación:

\begin{itemize}
    \item \texttt{THE\_MAN.py}: interfaz gráfica (GUI) basada en
          Tkinter que centraliza la entrada de parámetros, gestiona el
          proyecto y orquesta la ejecución y el postproceso. La
          ejecución del binario se lanza en un hilo aparte para que la
          interfaz no se bloquee, e incluye una vista previa del
          agregado de esferas.
    \item \texttt{mstm\_utils.py}: utilidades físicas; conversión a
          unidades adimensionales e interpolación del índice de
          refracción a la longitud de onda deseada.
    \item Scripts de representación (\texttt{plot\_nearfield.py},
          \texttt{plot\_scattering\_matrix.py},
          \texttt{plot\_radius\_sweep.py} y
          \texttt{plot\_asymmetry.py}): generan las figuras de campo
          cercano, campo lejano, barridos y parámetro de asimetría. La
          GUI los lanza como procesos independientes.
\end{itemize}

La comunicación entre Python y Fortran se realiza estrictamente por
archivos: la GUI escribe \texttt{mstm.inp}, lanza \texttt{mstm.exe}
mediante \texttt{subprocess} y recoge los \texttt{.dat} resultantes.
Cada simulación queda así autocontenida en una carpeta con su fichero
de entrada, sus salidas y su configuración serializada en JSON, lo que
garantiza la reproducibilidad de los resultados.

% ---------------------------------------------------------------------
\subsection{Parámetros físicos y unidades adimensionales}
\label{subsec:mstm_utils}

MSTM trabaja con todas las longitudes escaladas por el número de onda
en el medio externo,
\begin{equation}
    k_{0} = \frac{2\pi}{\lambda},
    \label{eq:length_scale}
\end{equation}
de modo que cada esfera queda caracterizada por su parámetro de tamaño
$x = k_{0}\,a$, con $a$ el radio, en línea con la convención de la
sección~\ref{sec:miecoeffs}. El módulo \texttt{mstm\_utils.py}
convierte los parámetros geométricos del usuario (longitud de onda y
radio en $\mu$m) al par (\texttt{length\_scale\_factor},
\texttt{size\_parameter}) que requiere el fichero de entrada.

A pesar de utilizar únicamente oro en este trabajo, 
el programa está configurado para
otros materiales seleccionables (Ag, Al, Cu,
H$_2$O). El módulo carga las tablas de \texttt{refractiveindex.info}
almacenadas localmente (Johnson \& Christy para los metales nobles,
Rakic para el aluminio, Hale \& Querry para el agua) e interpola
linealmente el índice complejo $\tilde{n} = n + i\,\kappa$ a la longitud de
onda solicitada. El espaciado de estas tablas ($\sim 10$\,nm) es
suficientemente fino para que el error de interpolación sea
despreciable frente a la incertidumbre experimental de los datos.
Cabe señalar que MSTM adopta el convenio de parte imaginaria positiva
para medios absorbentes, convención que el módulo respeta al leer las
tablas.

% ---------------------------------------------------------------------
\subsection{Barridos paramétricos}
\label{subsec:sweeps}

El entorno permite repetir una misma configuración variando un
parámetro, ya sea la longitud de onda o el radio de las esferas. En
lugar de invocar el binario una vez por cada valor, la GUI genera un
único \texttt{mstm.inp} con bloques \texttt{new\_run} sucesivos,
aprovechando el mecanismo nativo de MSTM. Los dos barridos son
excluyentes: solo uno puede estar activo a la vez.

En un barrido en longitud de onda, cada \texttt{new\_run} actualiza el
factor de escala $k_{0} = 2\pi/\lambda$ y el índice de refracción
$\tilde{n}(\lambda)$ interpolado, manteniendo fija la geometría física
del agregado. El resultado son las eficiencias y demás observables en
función de $\lambda$.

En un barrido en radio, la longitud de onda es fija, de modo que
$k_{0}$ y $\tilde{n}$ son constantes y se calculan una sola vez; lo
único que cambia entre \emph{runs} es el radio de las esferas. Cuando
se usa una de las geometrías predefinidas, los centros se recalculan en
cada paso con una separación centro a centro $d = 2r + s$, de manera
que la separación entre superficies $s$ se mantiene constante mientras
crece el radio. Esto permite estudiar cómo cambia la respuesta del
agregado con el tamaño de las partículas sin alterar el hueco entre
ellas.

Para facilitar este tipo de estudios, la GUI puede generar de forma
automática las posiciones de varias geometrías sencillas (dos o cuatro
esferas en línea, en cuadrado o en cuadrado girado $45^\circ$) y admite
asignar un radio común a todas las esferas o un radio distinto a cada
una.

% ---------------------------------------------------------------------
\subsection{Muestreo del campo cercano}
\label{subsec:adaptive_step}

El paso espacial $\Delta$ de la malla del campo cercano debe resolver
simultáneamente la longitud de onda, el radio de la esfera más pequeña
y el \emph{gap} entre esferas adyacentes. El módulo calcula
automáticamente un paso candidato como el mínimo de tres criterios:

\begin{equation}
    \Delta = \min\!\left(
        \frac{2\pi}{N_{\text{ppw}}},\;
        \frac{r_{\min}}{N_{\text{rad}}},\;
        \frac{g_{\min}}{N_{\text{gap}}}
    \right),
    \label{eq:dmin}
\end{equation}

con valores por defecto $N_{\text{ppw}}=50$, $N_{\text{rad}}=25$ y
$N_{\text{gap}}=12$, donde $r_{\min}$ es el radio mínimo del agregado
y $g_{\min} = \min_{i\neq j}\!\bigl(\|\mathbf{r}^i - \mathbf{r}^j\| -
a^i - a^j\bigr)$ es la separación superficial mínima entre pares de
esferas. Si la malla resultante superase $2\times 10^{6}$ puntos, el
paso se reescala hasta satisfacer esa cota. El usuario puede en
cualquier momento sustituir el valor automático por uno manual.

% ---------------------------------------------------------------------
\subsection{Visualización de resultados}
\label{subsec:visualizacion}

\subsubsection*{Campo cercano}

El script \texttt{plot\_nearfield.py} lee el fichero de campo cercano
de MSTM, que contiene para cada punto de la malla las componentes
complejas de los campos eléctrico y magnético correspondientes a las
dos polarizaciones de la onda incidente. A partir de ellas se calculan
la intensidad $|\mathbf{E}|^2$, el vector de Poynting promedio
$\langle\mathbf{S}\rangle = \frac{1}{2}\,\mathrm{Re}
[\mathbf{E}\times\mathbf{H}^{*}]$, y su promedio incoherente entre
polarizaciones para luz no polarizada,
$\frac{1}{2}(|\mathbf{E}_{\parallel}|^2 +
|\mathbf{E}_{\perp}|^2)$. Los resultados se representan como mapas de
contorno sobre la malla reconstruida, con las secciones transversales
de las esferas superpuestas.

\subsubsection*{Campo lejano}

El script \texttt{plot\_scattering\_matrix.py} extrae del fichero de
salida la matriz de Mueller y las eficiencias de extinción, absorción y
dispersión de cada \emph{run}. Representa la función de fase
$S_{11}(\theta)$ en coordenadas polares (en escalas lineal y
logarítmica de forma simultánea) y las eficiencias $Q_{\text{ext}}$,
$Q_{\text{abs}}$ y $Q_{\text{sca}}$ frente a la longitud de onda cuando
hay barrido espectral, o frente al radio cuando el barrido es en tamaño
(\texttt{plot\_radius\_sweep.py}).

El programa calcula también los grados de polarización del campo
dispersado a partir de sus parámetros de Stokes
(ec.~\ref{eq:MuellerMatrix}) y el parámetro de asimetría $g$
(ec.~\ref{eq:g_mie}, mediante \texttt{plot\_asymmetry.py}). Son
magnitudes disponibles en el entorno que se han empleado de forma
puntual según el caso.

\begin{figure}[H]
    \centering
    \includegraphics[width=1\textwidth]{Imagenes/Pantallazo_GUI.png}
    \caption{Interfaz gráfica del \emph{wrapper} en Python para MSTM.}
    \label{GUI:screenshot}
\end{figure}

En la figura~\ref{GUI:screenshot} se muestra una captura de la GUI con
la configuración de una simulación: cuatro esferas de oro, de radio
0.1\,$\mu$m, dispuestas en un cuadrado a 45$^\circ$ respecto a la onda
incidente. Se puede ver una vista previa del agregado, que se actualiza
en tiempo real a medida que se modifica la posición de las esferas, su
radio o la separación entre ellas. También se aprecia el tipo de
simulación que se va a ejecutar: un barrido en longitud de onda entre
400 y 800\,nm con 30 pasos, una onda plana incidente (el programa
también admite un haz gaussiano), con el orden multipolar determinado
automáticamente por el propio MSTM y sin cálculo del campo cercano.

Una vez hecha la configuración, los botones \texttt{Write mstm.inp} y
\texttt{Run MSTM Simulation} se encargan, respectivamente, de generar el
fichero de entrada con la sintaxis adecuada y de lanzar la simulación.
Cuando esta finaliza, los distintos botones \texttt{Plot} permiten
visualizar los resultados de forma inmediata.

Si el usuario carga una configuración guardada previamente, la GUI
reconstruye la geometría y los parámetros de la simulación a partir del
fichero JSON asociado y queda lista para lanzar la simulación o
visualizar los resultados, si estos ya están en la carpeta de la
simulación. Todos los resultados de este trabajo se han generado 
mediante esta interfaz.